\section{The half-residue extension for the actual pure-cubic labels}
\label{sec:pure-cubic-half-window}

\paragraph{Status and dependency boundary.}
This conditional-interface argument uses the original-coefficient
strict theorem AC-G154 for all three labels $1,\chi,\chi^2$.
It retains every omission already proved in AC-G153-window.
The class and local inputs used to prove AC-G154 have their own
recorded scopes; no new class computation is asserted here.
The remaining all-prime index estimate and mixed realization
are not supplied by this argument.

Put $p=269$, $K=\mathbb Q(\omega)$, $t^3=p$ with $t$ real,
and let $\chi:G_K\to\mu_3$ describe $K(t)/K$.
Let $\phi$ be the primitive odd character modulo sixteen with
$\phi(13)=i$ and kernel $\{1,7\}$, and put
$\psi_4=\chi_{-4}|_{G_K}$, $\xi_\pm=\chi_{\pm8}|_{G_K}$.
For a fixed source datum, $0<\varepsilon<1/256$, an allowed
page and its prescribed pair $q,d$, keep the actual fields
\[
 \begin{aligned}
 F_u&=\mathbb Q(\zeta_{2^{h(u)+2}},\zeta_7,
                           \sqrt{-q},\sqrt{-d},\omega),\\
 H_u&=F_u(t),\qquad G_u=\operatorname{Gal}(H_u/K),\quad D_u=|G_u|,\\
 u^{2+8\varepsilon}&\leq2^{h(u)}\leq2u^{2+8\varepsilon},
                      \qquad T=\lfloor u^{399}\rfloor.
 \end{aligned}
 \tag{PCW.1}
\]
All original coefficient rings $O_\ell$, fields $E_\ell$,
normalizing ranks $b_\ell$, full local conditions, derived
indices and the unrestricted coefficient-prime range are
exactly (FC.1)--(FC.2), (EU.2) and (TC.1).

For $u\geq2$ define the following disjoint sets of labels:
\[
 \begin{aligned}
 Z_{16}^+&=\{(p,\chi^{-j}\phi^{-\sigma},e):j=1,2,\ \sigma=\pm1,\\
 &\hspace{2em}1\leq e<T,\ e\equiv-1\pmod{67}\},\\
 Z_4^+&=\{(p,\chi^{-j}\psi_4,e):j=1,2,
       \ 1\leq e<T,\ e\equiv-1\pmod{67}\},\\
 Z_8&=\{(p,\chi^{-j}\xi_\nu,e):j=1,2,\ \nu=\pm,\\
 &\hspace{2em}1\leq e<T,\ e\equiv-1\pmod{134}\},\\
 Z_{\rm cub}^+&=\{(p,\chi^j,e):j=0,1,2,\ 1\leq e<T,\\
 &\hspace{2em}e\equiv0\pmod{268}\ \text{or}\ e\equiv-1\pmod{134}\},\\
 Z_3&=\{(3,1,e):1\leq e<T\},\\
 Z^{\rm cub+}&=Z_{16}^+\sqcup Z_4^+\sqcup Z_8
                                \sqcup Z_{\rm cub}^+\sqcup Z_3.
 \end{aligned}
 \tag{PCW.2}
\]
The former pure-cubic set in AC-G153 has the two residue classes
$0,267\pmod{268}$.  Hence the precise added set is
\[
 Z_{\rm cub}^{\rm new}=\{(p,\chi^j,e):j=0,1,2,
             \ 1\leq e<T,\ e\equiv133\pmod{268}\}.
 \tag{PCW.3}
\]
Let $I_{399}^{\rm estcub}$ be (TC.1) with exactly
$Z^{\rm cub+}$ omitted; the old remaining sum is
$I_{399}^{\rm est4+}$ from AC-G153.

\begin{theorem}[AC-G155-window: the pure-cubic half-residue reduction]
\label{AC-G155-window}
All labels in (PCW.2) occur on the prescribed host, and their
strict groups on the original coefficient modules vanish.
For positive integer $T$ the exact disjoint counts are
\[
 \begin{aligned}
 |Z_{16}^+|&=4\lfloor T/67\rfloor,&
 |Z_4^+|&=2\lfloor T/67\rfloor,\\
 |Z_8|&=4\lfloor T/134\rfloor,& |Z_3|&=T-1,\\
 |Z_{\rm cub}^+|&=3\bigl(\lfloor(T-1)/268\rfloor
                                  +\lfloor T/134\rfloor\bigr),\\
 |Z_{\rm cub}^{\rm new}|&=3\lfloor(T+134)/268\rfloor.
 \end{aligned}
 \tag{PCW.4}
\]
Let $C_u$ be the sum of the actual nonnegative normalized
positive-index costs over the full union $Z^{\rm cub+}$.  Then
\[
 \begin{aligned}
 I_{399}&=I_{399}^{\rm estcub}+C_u,\qquad
 I_{399}^{\rm estcub}\leq I_{399}^{\rm est4+}\leq I_{399},\\
 0\leq C_u&\leq8D_u^2T^2\log(8TD_u)
             +\tfrac12D_u\mathcal R_uT^2+2D_ur_uc_0T\\
 &=O_\varepsilon\bigl(u^{802+16\varepsilon}(\log u)^2\bigr).
 \end{aligned}
 \tag{PCW.5}
\]
The equality is also an identity of nonnegative extended sums,
without subtracting infinite quantities.  One full (FL.4)
majorant covers the entire old and new union.

The common-window premise (TC.3), or its AC-G153 form, may use
$I_{399}^{\rm estcub}$ with every simultaneous field, coefficient,
finite-character and Tate-weight quantifier unchanged.
Either old premise implies the new one with unchanged parameters.
The new premise implies the full old one using the same $W,J$,
source datum, page and pair, after increasing only its initial
threshold and halving the positive relative slack.
The actual source $R.t$, every mixed theta/Selmer and analytic
clause of $\mathcal R_{\rm cub}$, and all other AC-K35 inputs
remain required.  No such unproved input is asserted here.
\end{theorem}

\begin{proof}
\proofstep{1. The same three finite labels, including the trivial one.}
By definition $H_u$ contains $K(t)$, so $1,\chi,\chi^2$
already factor through $G_u$.
The field $F_u/\mathbb Q$ is abelian and cannot contain the
nonabelian $S_3$ field $K(t)$.
Since $K(t)/K$ has prime degree three, $F_u\cap K(t)=K$.
Its cubic subgroup makes these three characters distinct.
The other coefficient-$269$ families in (PCW.2) have nontrivial
quadratic or quartic two-power factors; the corresponding factors
remain nontrivial on $G_K$ because $K\cap\mathbb Q(\zeta_{16})
=\mathbb Q$.  Their order is prime to three and cannot be
canceled by a cubic factor.
The coefficient-three family has a different coefficient prime.
Thus all disjointness and old occurrence assertions of AC-G153
are retained, with no new finite character adjoined.

The original $O_p$ contains the cubic values because
$\exp(G_u)=3\cdot2^{h(u)}$.
For $j=0,1,2$ the new coefficient is literally
\[
 \begin{aligned}
 A_{j,e}&=(O_p\otimes_{\mathbb Z_p}\mathbb Q_p/\mathbb Z_p)
                                                  (\chi^j)(-e)\\
        &=(E_p/O_p)(\chi^jc_p^{-e}).
 \end{aligned}
 \tag{PCW.6}
\]
Apply (PCH.2) of AC-G154 with $m=0$, $\kappa=1$ and
its cubic index equal to $-j$ modulo three.  This gives exactly
(PCW.6), since negation permutes $\{0,1,2\}$.
The theorem directly includes $j=0$ on the original base $K$;
that case is not inferred from the nontrivial cubic sectors.
It permits the actual original $O_p$ with arbitrary absolute
ramification and gives strict zero for every
$e\equiv133\pmod{268}$.
The quadratic character at $269$ used in residual Kummer descent
does not occur in (PCW.6) or enlarge the host-label set.
Every old zero statement, including the coefficient-three result
without averaging by three, is retained from AC-G153.

\proofstep{2. Full local conditions and the exact additional count.}
The strict theorem uses zero in full
$H^1(G_{K_v},A_{j,e})$ at every $v\mid p$, and zero in
$H^1(I_v,A_{j,e})$ at each other finite place.
These are exactly the local conditions in (FC.2).
The coefficient (PCW.6) is unramified away from $3,p$,
including the case $j=0$.
At an added away-$p$ prime in $\Sigma_p$, inertia acts trivially.
A strict restriction that is a coboundary is therefore zero.
The cocycle identity makes the restriction zero on every conjugate
inertia subgroup and on the closed normal subgroup they generate.
The action and cocycle descend to the smaller ramification set;
conversely inflation is injective and preserves every remaining
local condition.  In particular no full local condition at $p$
is replaced by an unramified or weaker condition.

The new weights are $e=133+268n$, $n\geq0$.
For positive integer $T$, the number with $1\leq e<T$ is
$\lfloor(T+134)/268\rfloor$.
The two odd residue classes $133,267\bmod268$ are exactly
$e=134n-1$, $n\geq1$, counted by $\lfloor T/134\rfloor$.
The disjoint positive multiples of $268$ below $T$ are counted
by $\lfloor(T-1)/268\rfloor$.
Multiplying by the three distinct cubic labels proves (PCW.4),
with the useful identity
\[
 \lfloor T/134\rfloor=\lfloor T/268\rfloor
                               +\lfloor(T+134)/268\rfloor.
 \tag{PCW.7}
\]
All other counts are unchanged.  Precisely these strict terms
in (FC.15) vanish, without any assertion of individual index-zero.

\proofstep{3. A single full-union comparison bound.}
For each removed $(\ell,\eta,e)$ take the original dual complex
$C_{\ell,\eta,e}=R\Gamma(G_{K,\Sigma_\ell},
O_\ell(\eta^{-1})(e+1))$, with the codifferent convention (FL.2).
The determinant inequality (EU.2) gives
\[
 \begin{aligned}
 \frac{\log\#k_{O_\ell}}{b_\ell}\max(0,i_{\ell,\eta,e})
 &\leq\frac{\log\#H^2(C_{\ell,\eta,e})}{b_\ell}\\
 &=\frac{|\log\#H^2(C_{\ell,\eta,e})
                    -\log\#\mathcal S_{\ell,\eta}(e)|}{b_\ell}.
 \end{aligned}
 \tag{PCW.8}
\]
The equality uses the proved strict zero, retaining the actual
positive index cost and original normalizing rank $b_\ell$.
At a fixed weight, the entire union in (PCW.2) is one subsum
of the nonnegative all-prime error (FL.4), including the
coefficient-three terms.  Its cost is at most
\[
          16eD_u^2\log(8eD_u)+D_ue\mathcal R_u+2D_ur_uc_0.
 \tag{PCW.9}
\]
Summing once over $1\leq e<T$ and using $\sum e\leq T^2/2$,
$\log(8eD_u)\leq\log(8TD_u)$, $T-1\leq T$, proves (PCW.5).
No separate full error is charged for each family.
Nonnegativity gives additivity and monotonicity even when a
remaining unrestricted sum is infinite.

The same host bounds give $D_u\leq288u^{2+8\varepsilon}$,
$r_u\leq12$, and
$\mathcal R_u=O_\varepsilon(u^{\varepsilon/5}(\log u)^2)$,
uniformly over the prescribed page and pair.
With $T\leq u^{399}$ the two largest exponents are
$802+16\varepsilon$ and $800+41\varepsilon/5$.
Thus (PCW.5) is $o(\exp(u^\varepsilon/2))$ for each fixed
positive $\varepsilon$, with the same uniformity.

\proofstep{4. The same common window with its original quantifiers.}
Fix the datum, page and $0<\varepsilon<1/256$.
Suppose the new premise supplies $J\geq1$, $X_0$ and
$s\in(0,1)$, independent of $X,q,d,u,\ell,\eta,e$, so that
for every $X\geq X_0$ and prescribed
$\operatorname{WinPair}(\varepsilon,X^{1/J},q,d)$ it supplies
one $W$ with $X^{1/J}\leq W$, $W^2\leq X$ and
\[
 I_{399}^{\rm estcub}(H_u,u)\leq(1-s)\exp(u^\varepsilon/2)
                                 \quad(W\leq u\leq W^2).
 \tag{PCW.10}
\]
Choose $U_0\geq2$ so that (PCW.5) is at most
$(s/2)\exp(u^\varepsilon/2)$ for every $u\geq U_0$.
Only the retained datum, page, $\varepsilon$ and slack enter
this choice, not a new pair, coefficient, finite label or weight.
Replace $X_0$ by $\max(X_0,U_0^J)$.  The same $W$ gives
\[
 I_{399}(H_u,u)\leq(1-s/2)\exp(u^\varepsilon/2)
                                 \quad(W\leq u\leq W^2).
 \tag{PCW.11}
\]
For a fixed finite simultaneous family take the maximum of its
uniform error constants and keep this one $W$; no intersection
of separately selected windows is used.
Conversely either the full old premise or AC-G153's premise
implies (PCW.10) by nonnegativity with unchanged parameters.
The uniform little-$o$ formulations agree as well.
The actual host, rank $R.t$, family construction, mixed
theta/Selmer identification and analytic realization input
$\mathcal R_{\rm cub}$ remain exactly as before.
\end{proof}

\paragraph{The exact old-gate change.}
After AC-G153, EB-Q02 still includes $Z_{\rm cub}^{\rm new}$.
Its count is positive for $T\geq134$, so the new unknown sum
omits a proper additional subset of actual terms.
This is not a claim that the numerical sum strictly decreases.
Every previous omission is retained without new credit.
The remaining all-prime upper bound, mixed realization,
analytic comparison and H04.4 remain open.
The systems are the original rank-one Artin--Tate coefficients
on the same arithmetic curve and finite cover, with only the
previously audited SGA~1, Expos\'e~V, \S\S4 and 7 dictionary.
No new motive, comparison theorem or Lean execution is claimed.
